\documentclass[a4paper, 14pt]{extarticle}

\usepackage{cmap}
\usepackage[T2A]{fontenc}
\usepackage[utf8]{inputenc}
\usepackage[english, russian]{babel}
\usepackage{amsmath}
\usepackage{subcaption}
\usepackage{graphicx}
\usepackage{float}
\usepackage{hyperref}
\usepackage{multirow}

\usepackage[left=3cm, right=1.5cm, top=2cm, bottom=2cm]{geometry}
\usepackage{setspace}
\onehalfspacing 

\usepackage{indentfirst}
\setlength{\parindent}{1.25cm} 
\setlength{\emergencystretch}{3em}
\usepackage[expansion=false,protrusion=false,nopatch=footnote]{microtype}

\usepackage{tocloft}
\addto\captionsrussian{\renewcommand{\contentsname}{СОДЕРЖАНИЕ}}
\renewcommand{\cfttoctitlefont}{\hfill\normalfont}
\renewcommand{\cftaftertoctitle}{\hfill\vspace{\baselineskip}}
\renewcommand{\cftsecdotsep}{\cftdotsep}

\begin{document}

\begin{titlepage}
    \begin{center}
        Министерство науки и высшего образования Российской Федерации\\
        Федеральное государственное автономное образовательное учреждение высшего образования\\
        «Санкт-Петербургский политехнический университет Петра Великого»\\
        
        \textbf{Физико-механический институт}\\
        \textbf{Высшая школа прикладной математики и вычислительной физики}
    \end{center}

    \vfill
    
    \begin{center}
        \textbf{Отчёт по лабораторной работе \textnumero~ 7\\[0.5cm]}
        \textbf{Проверка гипотез о законе распределения. Критерий $\chi^2$\\}
        по дисциплине <<Математическая статистика>>
    \end{center}

    \vfill

    \hfill
    \begin{minipage}{0.45\textwidth}
        \textbf{Выполнил студент гр. 5030102/30202:}\\
        Шильдт Д.С.\\[0.5cm]
        \textbf{Преподаватель:}\\ 
        Баженов А.Н.
    \end{minipage}

    \vfill

    \begin{center}
        Санкт-Петербург\\2026
    \end{center}
\end{titlepage}

\setcounter{page}{2}

\tableofcontents
\newpage

\section{Постановка задачи}

\textbf{Задание:}
\begin{enumerate}
    \item Сгенерировать выборку объемом $n=100$ из $N(0,1)$.
    \item Оценить параметры $\mu,\sigma$ методом максимального правдоподобия (ММП). 
    \item Проверить гипотезу $H_0$: выборка из $N(\widehat{\mu},\widehat{\sigma})$ с помощью критерия $\chi^2$ при $\alpha=0.05$. 
    \item Исследовать чувствительность критерия: сгенерировать выборки из равномерного распределения и распределения Лапласа ($n=20$) и проверить их на нормальность.
\end{enumerate}

\textbf{Цели и задачи:}
\begin{itemize}
    \item Исследовать принцип работы критерия $\chi^2$ и его ограничения.
    \item  Оценить мощность критерия при альтернативных распределениях.
\end{itemize}

\section{Теоретическая часть}

\subsection{Используемые распределения}

Анализ чувствительности \textit{критерия согласия} осуществляется на базе трех непрерывных распределений с заданными плотностями вероятности:

\begin{enumerate}
    \item \textbf{Нормальный закон} $N(0, 1)$:
        \begin{equation*}
            f(x) = \frac{1}{\sqrt{2\pi}} \exp\left(-\frac{x^2}{2}\right)
        \end{equation*}
    \item \textbf{Распределение Лапласа} $L(0, 1/\sqrt{2})$:
        \begin{equation*}
            f(x) = \frac{1}{\sqrt{2}} \exp\left(-\sqrt{2}|x|\right)
        \end{equation*}
    \item \textbf{Равномерное распределение} $U(-\sqrt{3}, \sqrt{3})$:
        \begin{equation*}
            f(x) = \begin{cases} 
            \frac{1}{2\sqrt{3}}, & x \in [-\sqrt{3}, \sqrt{3}] \\ 
            0, & x \notin [-\sqrt{3}, \sqrt{3}] 
            \end{cases}
        \end{equation*}
\end{enumerate}

\subsection{Метод максимального правдоподобия}

Для оценки неизвестных параметров сдвига $m$ и масштаба $\sigma$ нормальной генеральной совокупности применяется метод максимального правдоподобия [1]. Функция правдоподобия $L$ для выборки $(x_1, \dots, x_n)$:
\begin{equation*}
\begin{aligned}
L(x_1, \dots, x_n, \mu, \sigma) &= \prod_{i=1}^{n} \frac{1}{\sigma \sqrt{2\pi}} \exp \left\{ -\frac{(x_i - \mu)^2}{2\sigma^2} \right\} = \\
&= (2\pi\sigma^2)^{-n/2} \exp \left\{ -\frac{1}{2\sigma^2} \sum_{i=1}^{n} (x_i - \mu)^2 \right\},
\end{aligned}
\end{equation*}

Оценки параметров определяются точкой глобального максимума данной функции. Переход к логарифмической функции правдоподобия упрощает дифференцирование:
\begin{equation*}
    \ln L = -\frac{n}{2}\ln(2\pi\sigma^2) - \frac{1}{2\sigma^2}\sum_{i=1}^n(x_i - \mu)^2.
\end{equation*}

Необходимое условие экстремума требует равенства нулю частных производных по оцениваемым параметрам $\mu$ и $\sigma^2$. Получим следующие уравнения правдоподобия:
\begin{equation*}
\left\{
\begin{aligned}
\frac{\partial \ln L}{\partial \mu} &= \frac{1}{\widehat{\sigma}^2} \sum_{i=1}^{n} (x_i - \widehat{\mu}) = \frac{n}{\widehat{\sigma}^2} (\bar{x} - \widehat{\mu}) = 0, \\
\frac{\partial \ln L}{\partial (\sigma^2)} &= -\frac{n}{2\widehat{\sigma}^2} + \frac{1}{2(\widehat{\sigma}^2)^2} \sum_{i=1}^{n} (x_i - \widehat{\mu})^2 = \\
&= \frac{n}{2(\widehat{\sigma}^2)^2} \left[ \frac{1}{n} \sum_{i=1}^{n} (x_i - \widehat{\mu})^2 - \widehat{\sigma}^2 \right] = 0,
\end{aligned}
\right.
\end{equation*}

Решение полученной системы нормальных уравнений устанавливает выборочное среднее в качестве точечной оценки математического ожидания и выборочную дисперсию в качестве оценки генеральной дисперсии [2]:
\begin{equation*}
    \hat{\mu} = \bar{x} = \frac{1}{n}\sum_{i=1}^n x_i, \quad \hat{\sigma} = \sqrt{\frac{1}{n}\sum_{i=1}^n (x_i - \hat{\mu})^2}.
\end{equation*}

\subsection{Критерий хи-квадрат Пирсона}

Проверка нулевой гипотезы $H_0$ реализуется с помощью непараметрического критерия $\chi^2$ Пирсона. Выборочное пространство разбивается на $k$ непересекающихся интервалов. \textit{Статистика критерия хи-квадрат} К.Пирсона [1] имеет вид:
\begin{equation*}
    \chi^2 = \sum_{i=1}^{k} \frac{n}{p_i} \left( \frac{n_i}{n} - p_i \right)^2 = \sum_{i=1}^{k} \frac{(n_i - n p_i)^2}{n p_i}.
\end{equation*}
\begin{flushleft}
где \\
\quad $k \approx 1.72\sqrt[3]{n}$ — формула Террелла; \\
\quad $k \approx 1 + 3.3\lg{n}$ — формула Стерджеса.
\end{flushleft}
Теоретическая вероятность $p_i$ вычисляется интегрированием плотности нормального распределения с найденными параметрами $\hat{\mu}$ и $\hat{\sigma}$ на $i$-м интервале группировки. 

Сходимость распределения статистики к закону хи-квадрат обеспечивается жестким ограничением $np_i \ge 5$. Число степеней свободы предельного распределения корректируется по теореме Фишера как $c = k - 1 - l$, где $l = 2$ — количество параметров, оцененных по исследуемой выборке. Критическая область формируется правосторонней: если вычисленная статистика $\chi^2$ превышает квантиль $\chi^2_{1-\alpha}(c)$ при заданном уровне значимости $\alpha$, нулевая гипотеза $H_0$ отвергается.

\section{Реализация}

\subsection{Язык программирования и среда}
\begin{itemize}
    \item \textbf{Используемый язык:} Python 3.14.2
    \item \textbf{Используемая среда разработки:} Visual Studio Code
\end{itemize}

\subsection{Используемые библиотеки}
\begin{itemize}
    \item \textbf{numpy} --- применяется для генерации выборок, вычисления точечных оценок параметров методом максимального правдоподобия, а также для объединения эмпирических частот.
    \item \textbf{scipy.stats} --- используется для определения теоретических вероятностей попадания случайной величины в интервалы группировки и вычисления критических квантилей предельного распределения $\chi^2$.
    \item \textbf{matplotlib.pyplot} --- используется для построения гистограмм и визуализации итоговых зависимостей.
    \item \textbf{pathlib} --- используется для автоматического определения пути и сохранения итоговых файлов таблиц и графиков.
\end{itemize}

\subsection{Суть реализованных методов}

\begin{itemize}
    \item \texttt{generate\_normal\_sample(n, rng)}, \texttt{generate\_uniform\_sample(n, rng)} и \texttt{generate\_laplace\_sample(n, rng)} --- функции формирования независимых выборок заданного объема $n$ из нормального закона $N(0, 1)$, равномерного $U(-\sqrt{3}, \sqrt{3})$ и распределения Лапласа $L(0, 1/\sqrt{2})$ соответственно.
    \item \texttt{estimate\_normal\_parameters(sample)} --- алгоритм расчета оценок параметров сдвига $\hat{\mu}$ и масштаба $\hat{\sigma}$ гипотетического нормального распределения методом максимального правдоподобия на основе исследуемой выборки.
    \item \texttt{choose\_bin\_count(n)} --- определяет количество интервалов группировки $k$ по формуле Стерджеса.
    \item \texttt{merge\_small\_frequency\_bins(observed, expected, edges, min\_freq=5)} --- процедура объединения соседних интервалов при нарушении условия $np_i \ge 5$. Алгоритм приоритизирует объединение граничных интервалов, затем обрабатывает внутренние интервалы с малыми ожидаемыми частотами.
    \item \texttt{chi\_square\_normal\_test(sample, alpha)} --- вычисления критерия Пирсона с применением процедуры объединения интервалов. Формирует границы интервалов, рассчитывает эмпирические частоты $n_i$ и теоретически ожидаемые $np_i$, выполняет объединение при необходимости. Возвращает значение вычисленной статистики $\chi^2$, скорректированное число степеней свободы и логический результат проверки гипотезы $H_0$.
    \item \texttt{evaluate\_sensitivity(...)} --- реализует цикл из 1000 итераций генерации выборок малого объема ($n=20$) из альтернативных законов и формирует эмпирическую оценку чувствительности критерия.
    \item \texttt{build\_breakdown\_rows(...)} --- формирует построчные данные разбиений для итоговых таблиц.
    \item \texttt{save\_single\_result\_csv(...)}, \texttt{save\_breakdown\_csv(...)}, \texttt{save\_sensitivity\_csv(...)} и \texttt{save\_sample\_csv(...)} --- функции сохранения рассчитанных таблиц в формате CSV.
    \item \texttt{save\_figures(...)} и \texttt{save\_sensitivity\_figure(...)} --- функции построения и сохранения графиков.
\end{itemize}

\section{Результаты}

В таблице \ref{tab:chi2_summary} приведены итоговые значения оценок параметров $\hat{\mu}$ и $\hat{\sigma}$, вычисленные значения статистики $\chi^2_{\textit{B}}$, соответствующие критические уровни и решение по нулевой гипотезе для каждой исследуемой выборки.

\begin{table}[H]
    \centering
    \caption{Сводные результаты проверки гипотез о нормальности распределения}
    \label{tab:chi2_summary}
    \begin{tabular}{|l|c|c|c|c|c|c|}
        \hline
        \textbf{Распределение} & $\mathbf{n}$ & $\mathbf{\hat{\mu}}$ & $\mathbf{\hat{\sigma}}$ & $\mathbf{\chi^2_{\textit{B}}}$ & $\mathbf{\chi^2_{1-\alpha}(c)}$ & \textbf{Результат} \\ \hline
        Нормальное & 100 & $-0.0503$ & $0.7728$ & $2.9944$ & $7.8147$ & Принимается \\ \hline
        Равномерное & 20 & $0.1384$ & $0.8759$ & $0.6573$ & $3.8415$ & Принимается \\ \hline
        Лапласа & 20 & $-0.2938$ & $0.9327$ & $3.4335$ & $3.8415$ & Принимается \\ \hline
    \end{tabular}
\end{table}

\begin{table}[H]
    \centering
    \caption{Разбиение и расчет $\chi^2$ для нормального распределения ($n=100$, $k=6$, $c=3$)}
    \label{tab:breakdown_norm}
    \begin{tabular}{|c|l|c|c|c|c|c|}
        \hline
        $\mathbf{i}$ & \textbf{Интервал} $\mathbf{\Delta_i}$ & $\mathbf{n_i}$ & $\mathbf{p_i}$ & $\mathbf{np_i}$ & $\mathbf{n_i - np_i}$ & $\mathbf{\frac{(n_i - np_i)^2}{np_i}}$ \\ \hline
        1 & $(-\infty, -0.9279]$ & 13 & 0.1281 & 12.8056 & 0.1944  & 0.0030 \\ \hline
        2 & $(-0.9279, -0.4163]$ & 19 & 0.1898 & 18.9828 & 0.0172  & 0.0000 \\ \hline
        3 & $(-0.4163, 0.0953]$  & 21 & 0.2568 & 25.6826 & -4.6826 & 0.8538 \\ \hline
        4 & $(0.0953, 0.6069]$   & 24 & 0.2277 & 22.7733 & 1.2267  & 0.0661 \\ \hline
        5 & $(0.6069, 1.1185]$   & 18 & 0.1323 & 13.2338 & 4.7662  & 1.7165 \\ \hline
        6 & $(1.1185, +\infty)$ & 5  & 0.0652 & 6.5218  & -1.5218 & 0.3551 \\ \hline
        \multicolumn{2}{|c|}{\textbf{СУММА}} & \textbf{100} & \textbf{1.0000} & \textbf{100.0000} & \textbf{0.0000} & \textbf{2.9944} \\ \hline
    \end{tabular}
\end{table}

\begin{table}[H]
    \centering
    \caption{Разбиение и расчет $\chi^2$ для равномерного распределения ($n=20$, $k=3$, $c=0$)}
    \label{tab:breakdown_uniform}
    \begin{tabular}{|c|l|c|c|c|c|c|}
        \hline
        $\mathbf{i}$ & \textbf{Интервал} $\mathbf{\Delta_i}$ & $\mathbf{n_i}$ & $\mathbf{p_i}$ & $\mathbf{np_i}$ & $\mathbf{n_i - np_i}$ & $\mathbf{\frac{(n_i - np_i)^2}{np_i}}$ \\ \hline
        1 & $(-\infty, -0.1888]$ & 7 & 0.3544 & 7.0871 & -0.0871  & 0.0011 \\ \hline
        2 & $(-0.1888, 0.4159]$ & 4 & 0.2700 & 5.3988 & -1.3988 & 0.3624 \\ \hline
        3 & $(0.4159, +\infty)$ & 9 & 0.3757 & 7.5141 & 1.4859  & 0.2938 \\ \hline
        \multicolumn{2}{|c|}{\textbf{СУММА}} & \textbf{20} & \textbf{1.0000} & \textbf{20.0000} & \textbf{0.0000} & \textbf{0.6573} \\ \hline
    \end{tabular}
\end{table}

\begin{table}[H]
    \centering
    \caption{Разбиение и расчет $\chi^2$ для распределения Лапласа ($n=20$, $k=3$, $c=0$)}
    \label{tab:breakdown_laplace}
    \begin{tabular}{|c|l|c|c|c|c|c|}
        \hline
        $\mathbf{i}$ & \textbf{Интервал} $\mathbf{\Delta_i}$ & $\mathbf{n_i}$ & $\mathbf{p_i}$ & $\mathbf{np_i}$ & $\mathbf{n_i - np_i}$ & $\mathbf{\frac{(n_i - np_i)^2}{np_i}}$ \\ \hline
        1 & $(-\infty, -0.6056]$ & 7  & 0.1691 & 7.3816 & -0.3816 & 0.0197 \\ \hline
        2 & $(-0.6056, 0.1832]$  & 10 & 0.3264 & 6.5284 & 3.4716  & 1.8460 \\ \hline
        3 & $(0.1832, +\infty)$ & 3  & 0.5045 & 6.0899 & -3.0899 & 1.5678 \\ \hline
        \multicolumn{2}{|c|}{\textbf{СУММА}} & \textbf{20} & \textbf{1.0000} & \textbf{20.0000} & \textbf{0.0000} & \textbf{3.4335} \\ \hline
    \end{tabular}
\end{table}

В таблице \ref{tab:sensitivity_summary} приведены показатели чувствительности критерия по 1000 повторениям: доля отвержений гипотезы и среднее значение статистики $\chi^2_{\textit{B}}$ для каждого альтернативного распределения.

\begin{table}[H]
    \centering
    \caption{Эмпирическая чувствительность критерия $\chi^2$ к отклонениям от нормальности (1000 итераций, $n=20$)}
    \label{tab:sensitivity_summary}
    \resizebox{\textwidth}{!}{%
    \begin{tabular}{|l|c|c|c|}
        \hline
        \textbf{Истинное распределение} & $\mathbf{n}$ & \textbf{Доля отвергнутых гипотез} & \textbf{Среднее значение} $\mathbf{\chi^2_{\textit{B}}}$ \\ \hline
        Равномерное $U(-\sqrt{3}, \sqrt{3})$ & 20 & $0.041$ & $0.950$ \\ \hline
        Лапласа $L(0, 1/\sqrt{2})$ & 20 & $0.086$ & $1.282$ \\ \hline
    \end{tabular}
    }
\end{table}

Графики гистограммы плотностей приведенных частот с наложением аналитических кривых плотности гипотетического нормального распределения $N(\hat{\mu}, \hat{\sigma})$.

\begin{figure}[H]
    \centering
    \includegraphics[width=1\textwidth]{../figures/chi_square_samples.png}
    \caption{Гистограммы выборок и плотности нормального распределения $N(\hat{\mu}, \hat{\sigma})$}
    \label{fig:chi_square_samples}
\end{figure}

\section{Обсуждение}
При проверке истинной нулевой гипотезы на нормальной выборке ($n=100$) вычисленная статистика $\chi^2_{\textit{B}} = 2.9944$ меньше критической квантили $7.8147$. Гипотеза правомерно принимается. Формирование первоначально $k_0=8$ интервалов было произведено по формуле Стерджеса. Применение условия $np_i \ge 5$ для минимальных ожидаемых частот вынудило объединить крайние интервалы, в результате чего было получено $k=6$ интервалов с сохранением $c=3$ степеней свободы для предельного распределения.

Анализ альтернативных законов на малых выборках ($n=20$) выявил критическую роль этапа объединения интервалов. При использовании первоначального разбиения ($k_0=5$) множество интервалов имели $np_i < 5$, нарушая требования теоремы. После объединения граничных интервалов остались только $k=3$ основных интервала, все с ожидаемыми частотами $\ge 5.3$. С учетом оценки двух параметров число степеней свободы стало $c=0$, что в реализации формально трактуется как $c=1$ (минимально допустимое значение для распределения хи-квадрат).

Для равномерного закона вычисленная статистика $\chi^2_{\textit{B}} = 0.6573$ значительно ниже критического порога $3.8415$, что приводит к принятию гипотезы о нормальности несмотря на явное отличие распределения. Объединение широких интервалов привело к сглаживанию структуры эмпирического распределения, маскируя истинное отклонение от нормальности.

Распределение Лапласа с тяжелыми хвостами демонстрирует близкую ситуацию. Вычисленная статистика $\chi^2_{\textit{B}} = 3.4335$ лишь немного меньше критической границы $3.8415$, хотя и не достигает уровня отклонения. Несмотря на наличие значительной асимметрии в эмпирических частотах (10 наблюдений в центральном интервале против 7 и 3 в краях), объединение интервалов размывает эти различия.

Результаты статистического моделирования (1000 итераций) количественно демонстрируют неработоспособность критерия Пирсона при малых объёмах выборок. Доля отвергнутых гипотез (мощность критерия) для равномерного распределения составляет всего $0.041$, для распределения Лапласа — $0.086$. Средние значения генерируемых статистик ($0.950$ и $1.282$ соответственно) стабильно лежат значительно ниже критического порога $3.8415$. Критерий обладает крайне низкой чувствительностью на выборках $n \le 20$ и демонстрирует систематическое совершение ошибок второго рода.

Проведённый анализ иллюстрирует фундаментальное ограничение критерия $\chi^2$: требование минимальных ожидаемых частот $np_i \ge 5$ хотя и обеспечивает теоретическую сходимость к предельному распределению, но на практике приводит к чрезмерному объединению интервалов на малых выборках, что в свою очередь значительно снижает информативность критерия.

\section{Выводы}

При проверке нулевой гипотезы на выборке из нормального распределения ($n=100$) критерий $\chi^2$ Пирсона корректно определил генеральный закон с применением процедуры объединения интервалов. Вычисленная статистика $\chi^2_{\textit{B}} = 2.9944$ строго меньше критической квантили $\chi^2_{0.95}(3) = 7.8147$.

Анализ альтернативных законов на выборках малого объема ($n=20$) выявил две взаимосвязанные проблемы. Во-первых, требование $np_i \ge 5$ приводит к необходимости объединения интервалов, что уменьшает число степеней свободы с $c=2$ до $c=0$ (факт.~$c=1$). Во-вторых, даже после объединения критерий остается неработоспособным: для равномерного распределения $\chi^2_{\textit{B}} = 0.6573 \ll 3.8415$, для распределения Лапласа $\chi^2_{\textit{B}} = 3.4335 < 3.8415$.

Массовое статистическое моделирование (1000 итераций) подтвердило стойкое отсутствие мощности критерия при $n \le 20$. Доля ошибочно отклоненных нормальных гипотез составила всего $4.1\%$ для равномерного и $8.6\%$ для лапласова распределения. Это свидетельствует о том, что критерий не способен различить совершенно не нормальные распределения на малых выборках.

Полученные результаты подчеркивают критическое значение этапа объединения интервалов с малыми ожидаемыми частотами и его негативное влияние на чувствительность критерия. Для надежного применения критерия хи-квадрат требуются выборки значительно большего объема (рекомендуется $n \ge 50$) или использование альтернативных тестов согласия (Колмогорова—Смирнова, Андерсона—Дарлинга), не требующих группировки данных.
 
\refstepcounter{section}
\addcontentsline{toc}{section}{\thesection\ Список литературы}
\renewcommand{\refname}{\thesection\ Список литературы}

\begin{thebibliography}{9}

    \bibitem{babakhina}
    Бабахина С. А., Баженов А. Н. Практикум по математической статистике: учеб. пособие. — СПб. : Санкт-Петербургский политехнический университет Петра Великого, 2026. 

    \bibitem{maksimov}
    Максимов Ю. Д. Математика. Теория и практика по математической статистике. Конспект-справочник по теории вероятностей : учеб. пособие. — СПб. : Изд-во Политехн. ун-та, 2009. — 395 с. 

\end{thebibliography}

\section{Приложение}

Исходный код программ доступен в репозитории на GitHub: \\
\url{https://github.com/Reichenau/math-stats-labs}
 
\end{document}
